
\documentclass{article}
\usepackage{colortbl}
\usepackage{makecell}
\usepackage{multirow}
\usepackage{supertabular}

\begin{document}

\newcounter{utterance}

\centering \large Interaction Transcript for game `mmlu\_pro', experiment `mmlu\_pro\_economics', episode 222 with Qwen3.6{-}27B.
\vspace{24pt}

{ \footnotesize  \setcounter{utterance}{1}
\setlength{\tabcolsep}{0pt}
\begin{supertabular}{c@{$\;$}|p{.15\linewidth}@{}p{.15\linewidth}p{.15\linewidth}p{.15\linewidth}p{.15\linewidth}p{.15\linewidth}}
    \# & \multicolumn{2}{c}{Player} && \multicolumn{2}{c}{Game Master} \\
    \hline

    \theutterance \stepcounter{utterance}  
    & & & \multicolumn{4}{p{0.6\linewidth}}{
        \cellcolor[rgb]{0.9,0.9,0.9}{
            \makecell[{{p{\linewidth}}}]{
                \texttt{\tiny{[P1$\langle$GM]}}
                \texttt{The following are multiple choice questions (with answers) about economics. Answer only with "The answer is (X)" where X is the correct letter choice.} \\
\\ 
\texttt{Question:} \\
\texttt{For a stationary autoregressive process, shocks will} \\
\texttt{Options:} \\
\texttt{A. Never occur} \\
\texttt{B. Persist indefinitely} \\
\texttt{C. Cause a permanent change in the trend} \\
\texttt{D. Grow exponentially} \\
\texttt{E. Only affect the next period} \\
\texttt{F. Eventually die away} \\
\texttt{G. Cause the process to become non{-}stationary} \\
\texttt{H. Immediately disappear} \\
\texttt{I. Lead to a cyclical pattern} \\
\texttt{J. Be balanced out by subsequent shocks} \\
\texttt{The answer is (F)} \\
\\ 
\texttt{Question:} \\
\texttt{Consider the following AR(1) model with the disturbances having zero mean and unit variance} \\
\texttt{yt = 0.2 + 0.4 yt{-}1 + ut} \\
\texttt{The (unconditional) mean of y will be given by} \\
\texttt{Options:} \\
\texttt{A. 0.45} \\
\texttt{B. 0.2} \\
\texttt{C. 0.4} \\
\texttt{D. 0.1} \\
\texttt{E. 0.3} \\
\texttt{F. 0.55} \\
\texttt{G. 0.25} \\
\texttt{H. 0.33} \\
\texttt{I. 0.6} \\
\texttt{J. 0.5} \\
\texttt{The answer is (H)} \\
\\ 
\texttt{Question:} \\
\texttt{Suppose that a test statistic has associated with it a p{-}value of 0.08. Which one of the following statements is true?} \\
\texttt{(i) If the size of the test were exactly 8\%, we would be indifferent between rejecting and not rejecting the null hypothesis} \\
\texttt{(ii) The null would be rejected if a 10\% size of test were used} \\
\texttt{(iii) The null would not be rejected if a 1\% size of test were used} \\
\texttt{(iv) The null would be rejected if a 5\% size of test were used.} \\
\texttt{Options:} \\
\texttt{A. (iii) and (iv) only} \\
\texttt{B. (i) and (iv) only} \\
\texttt{C. (i), (ii), and (iii) only} \\
\texttt{D. (i) and (ii) only} \\
\texttt{E. (i) and (iii) only} \\
\texttt{F. (i), (ii), (iii), and (iv).} \\
\texttt{G. (ii), (iii), and (iv) only} \\
\texttt{H. (ii) and (iii) only} \\
\texttt{I. (ii) and (iv) only} \\
\texttt{J. (i), (iii), and (iv) only} \\
\texttt{The answer is (C)} \\
\\ 
\texttt{Question:} \\
\texttt{Which of the following policies best describes supply{-}side fiscal policy?} \\
\texttt{Options:} \\
\texttt{A. Higher taxes on household income} \\
\texttt{B. Increased government spending} \\
\texttt{C. Increased taxes on corporate profits} \\
\texttt{D. Increased import tariffs} \\
\texttt{E. Decreased interest rates} \\
\texttt{F. Lower taxes on consumer goods} \\
\texttt{G. Lower taxes on research and development of new technology} \\
\texttt{H. Reduced government spending} \\
\texttt{I. Higher taxes on capital gains} \\
\texttt{J. An increase in the money supply} \\
\texttt{The answer is (G)} \\
\\ 
\texttt{Question:} \\
\texttt{The concentration ratio for a monopoly is} \\
\texttt{Options:} \\
\texttt{A. 50} \\
\texttt{B. 5} \\
\texttt{C. 10} \\
\texttt{D. 90} \\
\texttt{E. 15} \\
\texttt{F. 100} \\
\texttt{G. 0} \\
\texttt{H. 25} \\
\texttt{I. 75} \\
\texttt{J. N/A} \\
\texttt{The answer is (F)} \\
\\ 
\texttt{Question:} \\
\texttt{Monopoly deadweight loss is the result of} \\
\texttt{Options:} \\
\texttt{A. monopoly output being less than the competitive output.} \\
\texttt{B. setting the price below average total cost.} \\
\texttt{C. short{-}run abnormal profits.} \\
\texttt{D. monopoly output being greater than the competitive output.} \\
\texttt{E. long{-}run abnormal profits.} \\
\texttt{F. setting the price below marginal cost.} \\
\texttt{G. long{-}run normal profits.} \\
\texttt{H. setting the price equal to average total cost.} \\
\texttt{I. setting the price above average total cost.} \\
\texttt{J. setting the price above marginal cost.} \\
            }
        }
    }
     \\ \\

    \theutterance \stepcounter{utterance}  
    & \multicolumn{4}{p{0.6\linewidth}}{
        \cellcolor[rgb]{0.561,0.737,0.561}{
            \makecell[{{p{\linewidth}}}]{
                \texttt{\tiny{[P1$\rangle$GM]}}
                \texttt{The answer is (A)} \\
            }
        }
    }
    & & \\ \\

    \theutterance \stepcounter{utterance}  
    & & & \multicolumn{2}{p{0.3\linewidth}}{
        \cellcolor[rgb]{0.9,0.9,0.9}{
            \makecell[{{p{\linewidth}}}]{
                \texttt{\tiny{[GM$|$GM]}}
                \texttt{A} \\
            }
        }
    }
    & & \\ \\

    \theutterance \stepcounter{utterance}  
    & & & \multicolumn{2}{p{0.3\linewidth}}{
        \cellcolor[rgb]{0.9,0.9,0.9}{
            \makecell[{{p{\linewidth}}}]{
                \texttt{\tiny{[GM$|$GM]}}
                \texttt{J} \\
            }
        }
    }
    & & \\ \\

    \theutterance \stepcounter{utterance}  
    & & & \multicolumn{2}{p{0.3\linewidth}}{
        \cellcolor[rgb]{0.9,0.9,0.9}{
            \makecell[{{p{\linewidth}}}]{
                \texttt{\tiny{[GM$|$GM]}}
                \texttt{game\_result = LOSE} \\
            }
        }
    }
    & & \\ \\

\end{supertabular}
}

\end{document}
